%% ============================================================
%%  Chapter 3 - System Architecture
%% ============================================================
\chapter{System Architecture}
\label{ch:arch}

\section{Architectural Boundary Conditions}
\label{sec:boundary}

The architecture is bounded by six deployment constraints:
\begin{enumerate}
  \item Every participating AP must be administratively enrolled.
  \item Subscriber identity must be stronger than a shared Wi-Fi password; certificate-based or profile-based \eap{} is the minimum proposed mechanism.
  \item The home domain remains the authoritative source of subscriber policy.
  \item Authentication, policy compliance, slice isolation, accounting consistency, and edge-resource feasibility are hard feasibility requirements.
  \item Commodity edge APs are treated as constrained compute platforms.
  \item Machine learning operates only in the slow control path and cannot override local safety rules.
\end{enumerate}

\section{Clean Four-Plane Architecture}
\label{sec:clean-architecture}

The architecture is organised into four planes: trust and policy, fast edge enforcement, backhaul and gateway selection, and telemetry/optimisation/survivability. Figure~\ref{fig:architecture-clean} presents the simplified structure used throughout the rest of the proposal.

%% ============================================================
%%  Clean Layered Architecture Figure (readable, no clipping)
%% ============================================================
\begin{figure}[H]
\centering
\begin{tikzpicture}[
  font=\small,
  layer/.style={rectangle, rounded corners, draw=black, thick, align=center, text width=11.2cm, minimum height=1.25cm},
  arrow/.style={-{Latex[length=2.8mm]}, thick},
  dasharrow/.style={-{Latex[length=2.8mm]}, thick, dashed}
]

\node[layer, fill=softblue] (trust) at (0,0) {\textbf{Trust and policy plane}\\Subscriber identity $\rightarrow$ AP authenticator $\rightarrow$ AAA/RadSec federation $\rightarrow$ home policy/accounting};
\node[layer, fill=softgreen] (edge) at (0,-1.95) {\textbf{Fast edge enforcement plane}\\Safety gate: $Auth\land Pol\land Iso\land Acct$\\VLAN/tc/HTB slicing $\rightarrow$ queue and VXLAN/WireGuard tunnel};
\node[layer, fill=softorange] (backhaul) at (0,-4.05) {\textbf{Backhaul and gateway plane}\\Backhaul selector $\rightarrow$ terrestrial path by default\\optional LEO/Starlink-class path $\rightarrow$ home/sovereign gateway};
\node[layer, fill=softgray] (analytics) at (0,-6.25) {\textbf{Telemetry, optimisation, and survivability plane}\\RF/queue/accounting telemetry $\rightarrow$ graph state $\bar{G}_{\tau}$ $\rightarrow$ slow GNN/CMDP update\\token cache and signed offline logs for degraded mode};

\draw[arrow] (trust.south) -- node[right]{policy attributes} (edge.north);
\draw[arrow] (edge.south) -- node[right]{safe session} (backhaul.north);
\draw[dasharrow] (backhaul.south) -- node[right]{telemetry} (analytics.north);
\draw[dasharrow] (analytics.west) -- ++(-0.8,0) |- node[pos=0.22,left]{slow feedback only} (trust.west);

\node[align=center, font=\footnotesize] at (0,-7.55) {\textbf{Control split:} fast path = local AP/gateway, $1$--$100$ ms; slow path = regional/controller, $10$--$60$ s.};

\end{tikzpicture}
\caption{Clean layered subscriber-centric broadband architecture. The packet fast path remains local at the edge AP, while the GNN/CMDP block only issues slow macro-policy updates. Optional LEO backhaul is isolated to the backhaul plane and cannot bypass the safety gate.}
\label{fig:architecture-clean}
\end{figure}


\paragraph{Trust and policy plane.} Handles federated AAA, \radsec{} transport, home-policy lookup, access decisions, and accounting identity initialisation.

\paragraph{Fast edge enforcement plane.} Executes packet forwarding, VLAN/VXLAN separation, Linux \texttt{tc}/\htb{}/\fqcodel{} shaping, local overload protection, and isolation checks.

\paragraph{Backhaul and gateway plane.} Selects terrestrial or optional LEO backhaul paths after authentication and slice assignment have succeeded.

\paragraph{Telemetry, optimisation, and survivability plane.} Aggregates telemetry into graph states, performs slow macro-policy optimisation, and supports token-cache degraded mode during local isolation.

\section{Fast Path, Slow Path, and Control Epochs}
\label{sec:timescales}

The two-timescale model prevents the GNN controller from chasing millisecond-scale packet bursts. The local AP and edge gateway enforce existing rules at fast timescales, while the GNN updates macro policies only after telemetry aggregation.

\begin{equation}
  \Tfast \ll \TGNN, \qquad t_{\mathrm{infer}} \leq 150\,\mathrm{ms} \ll \TGNN.
  \label{eq:timescale-sep}
\end{equation}

\begin{table}[H]
  \centering
  \caption{Control-plane split and temporal responsibilities.}
  \label{tab:control-split}
  \begin{tabularx}{\textwidth}{lllX}
    \toprule
    \textbf{Layer} & \textbf{Location} & \textbf{Timescale} & \textbf{Responsibility} \\
    \midrule
    Fast path & Edge AP/gateway & $1$--$100\,\mathrm{ms}$ & Enforce slice rules, queue shaping, VLAN/tunnel isolation, overload protection. \\
    Local guard & Edge AP/gateway & $100\,\mathrm{ms}$--$1\,\mathrm{s}$ & Apply deterministic fallback under overload, queue instability, or backhaul loss. \\
    GNN control & Regional controller & $10$--$60\,\mathrm{s}$ & Update association preferences, channel plans, slice weights, and backhaul policy. \\
    Offline trainer & Research server & Hours--days & Train and validate graph models from telemetry and simulation traces. \\
    \bottomrule
  \end{tabularx}
\end{table}

\section{Commodity Edge Resource Budget}
\label{sec:edge-budget}

Each edge AP is modelled as a constrained compute platform with resource vector
\begin{equation}
  \Redge_j=(CPU_j,RAM_j,IRQ_j,Crypto_j,Encaps_j,Queue_j,Tel_j).
  \label{eq:resource-vector}
\end{equation}

The feasibility constraint is
\begin{equation}
  CPU_j\leq CPU_j^{\max},\quad RAM_j\leq RAM_j^{\max},\quad IRQ_j\leq IRQ_j^{\max},\quad \forall j.
  \label{eq:edge-resource-ceiling}
\end{equation}

The operational acceptance target is $CPU_j\leq0.75CPU_j^{\max}$ under peak traffic, leaving headroom for control-plane tasks. This target is an empirical acceptance criterion, not a theoretical guarantee.

\subsection{Driver Buffer Admission Heuristic}
\label{sec:ring-buffer}

The driver-ring control is an empirical admission-control heuristic, not a stability theorem. Let $\Lambda_{k,j}$ be the measured normalised frame load of slice $k$ on AP $a_j$. The admission condition is
\begin{equation}
  \Omega_j=\sum_{k\in\mathcal{K}}\Lambda_{k,j}
  \leq
  \Omega_j^{\max}
  \left[1-\xi\frac{Drop_{driver,j}}{Arrival_j+\epsilon}\right]_+,
  \label{eq:ring-buffer}
\end{equation}
where $[x]_+=\max(0,x)$, $\epsilon>0$ prevents division by zero, and $\xi$ is selected through hardware calibration. This prevents the earlier draft from asserting an unproven universal ring-buffer stability condition.

\section{Edge Survivability Under Local Isolation}
\label{sec:survivability}

Nepal-specific power cuts, fibre disruptions, and intermittent backhaul motivate a three-state survivability model:
\begin{equation}
  \omega_j\in\{\mathrm{Online},\mathrm{Degraded},\mathrm{Isolated}\}.
\end{equation}

In \emph{Online} mode, all authentication and accounting flows use the normal AAA/RadSec path. In \emph{Degraded} mode, previously authorised subscribers may use short-lived cached tokens with strict bandwidth caps and signed offline accounting. In \emph{Isolated} mode, no privilege escalation is possible; the AP provides either no subscriber service or limited guest-degraded service according to local policy.

A cached token is
\begin{equation}
  \tau_i=\hmaclabel{}_{\Knet}(ID_i\|TS_i\|SLA_i),
  \label{eq:hmac-token}
\end{equation}
with anti-replay timestamp $TS_i$, service entitlement bitmap $SLA_i$, and key $\Knet{}$ provisioned before failure. The CPU target for re-authentication storms is treated as an empirical acceptance criterion:
\begin{equation}
  CPU_{auth,j}\leq0.05CPU_j^{\max}
  \quad \text{during a 100-client cached-token re-authentication test.}
  \label{eq:hmac-cpu-criterion}
\end{equation}

\section{Normalised LEO Control Backoff}
\label{sec:leo-backoff}

The clipping operator is defined as

\begin{equation}
\operatorname{clip}_{[a,b]}(x)
=
\min\{b,\max\{a,x\}\}.
\label{eq:clip-definition}
\end{equation}

The LEO-specific macro-epoch backoff uses dimensionless inputs:

\begin{equation}
\Delta t_{\mathrm{eff}}^{\mathrm{LEO}}
=
\operatorname{clip}_{[\Delta t_{\mathrm{macro}},\Delta t_{\max}]}
\left(
\Delta t_{\mathrm{macro}}
\exp
\left[
\kappa_{\mathrm{RTT}}
\frac{\sigma^2_{\mathrm{RTT}}}{\sigma^2_{\mathrm{RTT,ref}}+\epsilon}
+
\kappa_{\mathrm{drop}}
\frac{Drop_{\mathrm{LEO}}}{Drop_{\mathrm{ref}}+\epsilon}
\right]
\right).
\label{eq:leo-backoff-clean}
\end{equation}

The clipping operation prevents exponential blow-up, while normalisation
prevents dimensional inconsistency inside the exponential. The coefficients
$\kappa_{\mathrm{RTT}}\geq0$ and $\kappa_{\mathrm{drop}}\geq0$ are LEO-backoff
sensitivity parameters. They are not LEO cost weights and are therefore not
constrained to sum to one.

\section{Systems Verification Matrix}
\label{sec:verification}

\begin{table}[H]
  \centering
  \caption{Systems verification and boundary performance matrix.}
  \label{tab:verification}
  \begin{tabularx}{\textwidth}{llX}
    \toprule
    \textbf{Element} & \textbf{Environment} & \textbf{Acceptance Target} \\
    \midrule
    Slice isolation & Emulation + prototype & Zero cross-slice packets matching foreign VLAN/VNI tags. \\
    Edge compute ceiling & Hardware benchmark & $CPU_j\leq0.75CPU_j^{\max}$ under peak configured workload. \\
    Auth-storm resistance & Hardware benchmark & $CPU_{auth,j}\leq0.05CPU_j^{\max}$ as empirical target, not guarantee. \\
    GNN inference & Trace simulation & $t_{infer}\leq150\,\mathrm{ms}$ and $t_{infer}\ll \TGNN$. \\
    Wi-Fi path-loss model & RSSI survey & RMSE target $4$--$6\,\mathrm{dB}$ against floor-plan/RSSI heatmap data. \\
    LEO model & Backhaul tests/simulation & LEO results reported separately from terrestrial-only baseline. \\
    Accounting consistency & Prototype + emulation & $|Bytes_{RADIUS}-Bytes_{GW}|\leq\varepsilon_{acct}$. \\
    \bottomrule
  \end{tabularx}
\end{table}
